\section{Several observations in a logarithmic return window}
\label{sec:logarithmic-return-windows}
A product of observations at different actual returns is not covered
by assuming a variation bound for one of its factors. We prove that
such a product has a controlled approximation when its return window
grows logarithmically. The final event and the record remain unchanged;
the collision cutoff occurs only in the approximation used in the proof.

Let $0\le\ell_1<\cdots<\ell_q\le m$ and let $u_1,\ldots,u_q$
be complex section functions with $\|u_j\|_\infty\le1$ and
zero extensions $u_j^0\in\mathrm{BV}(M)$. Set
\begin{equation}\label{eq:window-weight-definition}
 W_R(y)=\prod_{j=1}^{q}u_j((F_R^*)^{\ell_j}y),\qquad
 V=1+\sum_{j=1}^{q}\|u_j^0\|_{\mathrm{BV}},\qquad
 \alpha_R=\int_{Y_R^*}W_R\dd\nu_R^*.
\end{equation}
All the data may depend on $n,R$. Empty products are one. No
factorization of $\alpha_R$ into individual observation masses is
asserted. For $0\le k\le n-m$ put
\begin{equation}\label{eq:window-transform}
 \mathcal C_{n,k,R}^{W}(v)=
 \int_{Y_R^*}W_R((F_R^*)^kx)
 e^{iv\cdot U_{n,R}(x)/\sqrt n}\dd\nu_R^*(x),
 \qquad U_{n,R}=J_{n,R}-n\bar G_R.
\end{equation}
This convention uses the normalized section probability. Thus all the
observations occur at their actual return states between returns $k$
and $k+m$, including a terminal window when $k=n-m$.

\subsection{A controlled approximation to the whole window}
\begin{lemma}[Truncated and smoothed window weights]
\label{lem:window-weight-approximation}
There is $C_{\mathrm w}\ge1$, uniform in $R,m,q,L$, such that for
integers $L\ge\max\{m,1\}$ and $0<\epsilon<1/2$ there is a section
function $W_{L,\epsilon}$ with
\begin{align}
 \|W_{L,\epsilon}\|_\infty&\le1,\qquad
 \|(W_{L,\epsilon})^0\|_{\mathrm{BV}}
       \le C_{\mathrm w}(1+q\epsilon^{-1})e^{C_{\mathrm w}L},
       \label{eq:window-weight-BV}\\
 \|W_R-W_{L,\epsilon}\|_{L^1(\nu_R^*)}
       &\le C\{\epsilon V+e^{a_{\mathrm t}m-c_{\mathrm t}L}\}.
       \label{eq:window-weight-L1}
\end{align}
The tail constants are those of
\eqref{eq:path-cumulative-tail-convention}.
\end{lemma}
\begin{proof}
Let $v_j=\mathsf S_\epsilon u_j^0$ and define, on the actual section,
\[
 W_{L,\epsilon}(y)=\one_{\{N_{m,R}(y)\le L\}}
                \prod_{j=1}^{q}v_j((F_R^*)^{\ell_j}y).
\]
Extend this function by zero to $M$. Its supremum is at most one.
Telescoping products and invariance of the actual first-return map give
\[
 \int_{Y_R^*}\left|
  \prod_j u_j\circ(F_R^*)^{\ell_j}
  -\prod_j v_j\circ(F_R^*)^{\ell_j}\right|\dd\nu_R^*
 \le \sum_j\|u_j-v_j\|_{L^1(\nu_R^*)}\le C\epsilon V.
\]
This use of invariance is valid before imposing the count cutoff.
The cutoff changes the product by at most its event probability, which
is bounded by \eqref{eq:path-cumulative-tail-convention}. This proves
\eqref{eq:window-weight-L1}.

To prove the variation assertion, resolve $\{N_{m,R}\le L\}$ by
its first $L$ binary section-membership symbols. There are at most
$2^L$ possibilities, with fixed initial membership equal to one. On
each such word the collision indices of the first $m$ actual returns
are fixed integers between zero and $L$. Hence the product above
is a product of $q$ smooth functions $v_j\circ T_R^{t_j}$, multiplied
by the word indicator and the initial section indicator. By
Theorem~\ref{thm:exponential-finite-record} and
$\|v_j\|_{C^1}\le C\epsilon^{-1}$, the BV product inequality bounds
its variation by $C(1+q\epsilon^{-1})e^{CL}$ on each word.
Summing the words and enlarging the exponential constant gives
\eqref{eq:window-weight-BV}. All internal word jumps and the initial
section trace are included. The summands have disjoint supports, which
is why their number does not enter the supremum bound. Null singular
sets are assigned the same convention as in the finite-record theorem.
\end{proof}

\subsection{A variable central band}
\begin{lemma}[The marked estimate on a specified subband]
\label{lem:variable-marked-band}
For $0\le\theta\le1/200$, define
\[
 s_0(\theta)=1/28-5\theta,\qquad
 s_1(\theta)=1/14-4\theta.
\]
For a bounded complex section function $a$, with zero-extension
variation $V_a$, supremum $M_a$, and mass
$\alpha_a=\int a\dd\nu_R^*$, the marked transform under $\nu_R^*$
satisfies
\begin{align}
 &\int_{|v|\le2n^\theta}
 \left|\int a((F_R^*)^kx)e^{iv\cdot U_{n,R}/\sqrt n}\dd\nu_R^*
                 -\alpha_a e^{-v^{\mathsf T}D_Rv/2}\right|\dd v
 \notag\\
 &\hspace{15mm}\le C\{M_a n^{-s_0(\theta)}\sqrt{\log(2+n)}
                           +V_a n^{-s_1(\theta)}\}.
       \label{eq:variable-marked-band}
\end{align}
The constants may be chosen on the full displayed interval of $\theta$,
and uniformly in $0\le k\le n$. The physical radius is
$2n^{-1/2+\theta}$, not $2n^\theta$.
\end{lemma}
\begin{proof}
Use the insertion $c_*^{-1}a^0$ in
Lemma~\ref{lem:marked-spectral} and retain
$\delta=\tfrac14n^{-1/14}$ and a logarithmic spectral block length.
Replace only the rescaled ball by $|v|\le2n^\theta$.
The observable-smoothing term integrates to
$C M_a n^{5\theta-1/28}\sqrt{\log(2+n)}$ and the insertion-smoothing
term to $CV_a n^{4\theta-1/14}$.
The remaining polynomial margins are
\[
 1/14-6\theta,\quad 3/14-5\theta,\quad 1/14-7\theta,
 \quad 1/5-4\theta,\quad 1/5-5\theta,
 \quad 1/2-5\theta,\quad 1-6\theta.
\]
They are strictly larger than $s_0(\theta)$ throughout
$0\le\theta\le1/200$. Their logarithmic factors are therefore
absorbed uniformly; the complementary power is smaller still. The two
analytic smallness conditions hold on the largest ball and hence on
all the smaller balls. This proves \eqref{eq:variable-marked-band}.
At $\theta=1/200$ it gives exactly the exponents $3/280$ and $9/175$
of Theorem C. Thus no inherited central estimate is replaced.
\end{proof}

\begin{theorem}[Central integrals with a full return window]
\label{thm:window-central-integral}
For the weight \eqref{eq:window-weight-definition}, $n\ge2$,
$0\le k\le n-m$, $L\ge\max\{m,1\}$, $0<\epsilon<1/2$, and
$0\le\theta\le1/200$,
\begin{align}
 &\int_{|v|\le2n^\theta}
 \left|\mathcal C_{n,k,R}^{W}(v)
                       -\alpha_R e^{-v^{\mathsf T}D_Rv/2}\right|\dd v
 \notag\\
 &\quad\le C\left\{
 n^{-s_0(\theta)}\sqrt{\log(2+n)}
 +(1+q\epsilon^{-1})e^{C_{\mathrm w}L}n^{-s_1(\theta)}
 +n^{4\theta}\bigl(\epsilon V+
                         e^{a_{\mathrm t}m-c_{\mathrm t}L}\bigr)
                      \right\}.
       \label{eq:window-central-integral}
\end{align}
The constant is independent of the location and the length of the
window and of its observation functions. The mass $\alpha_R$ and
the transform on the left use the full original window, with no
collision cutoff and no smoothing of the observations.
\end{theorem}
\begin{proof}
Apply Lemma~\ref{lem:variable-marked-band} to $W_{L,\epsilon}$ at
return $k$. Lemma~\ref{lem:window-weight-approximation} controls its
supremum and variation. Replacing that weight by $W_R$ changes the
transform pointwise by at most the $L^1$ error in
\eqref{eq:window-weight-L1}; invariance allows its evaluation at
$(F_R^*)^k x$. Integration over the four-dimensional ball costs
$Cn^{4\theta}$ times that error. The difference of the two masses
is bounded by the same error. Uniform positive definiteness gives
\[
 \int_{\R^4}e^{-v^{\mathsf T}D_Rv/2}\dd v\le C,
\]
so the Gaussian mass correction is no larger. This proves the result.
At no step is a randomly evaluated observation replaced by the same
observation at a deterministic collision time.
\end{proof}

\subsection{An unchanged multiple-time event}
\begin{corollary}[Logarithmic windows and rare-event budgets]
\label{cor:logarithmic-window-conditioning}
Suppose, for fixed constants $a,b,d>0$, $\kappa\ge0$, and $r\ge0$,
\[
 m\le a\log n,
 \qquad V\le C n^\kappa[\log(2+n)]^r,
 \qquad b>a.
\]
Take $L=\lceil b\log n\rceil$ and $\epsilon=n^{-d}/4$.
For $0<\theta\le1/200$ put
\begin{align}
 \sigma_1&=1/28-5\theta,&
 \sigma_2&=1/14-4\theta-C_{\mathrm w}b-d,\notag\\
 \sigma_3&=d-\kappa-4\theta,&
 \sigma_4&=c_{\mathrm t}b-a_{\mathrm t}a-4\theta.
       \label{eq:window-exponent-margins}
\end{align}
If all four numbers are positive, the integral in
\eqref{eq:window-central-integral} tends to zero at a power rate,
uniformly in $R,k$ and the permitted data.

For indicators $u_j=\one_{E_j}$, let
\begin{equation}\label{eq:exact-window-event}
 A_{n,k,R}=\bigcap_{j=1}^{q}
       \{x:(F_R^*)^{k+\ell_j}x\in E_j\},\qquad
 p_{n,R}=\nu_R^*(A_{n,k,R}).
\end{equation}
This probability is independent of $k$. If
$p_{n,R}\ge p_0n^{-\beta}>0$ and
$\beta<\min_i\sigma_i$, then
\begin{equation}\label{eq:conditioned-window-band}
 \int_{|v|\le2n^\theta}
 \left|\mathbb E_{\nu_R^*}
   [e^{iv\cdot U_{n,R}/\sqrt n}\mid A_{n,k,R}]
                      -e^{-v^{\mathsf T}D_Rv/2}\right|\dd v
                         \longrightarrow0
\end{equation}
uniformly. The event in this conclusion is exactly
\eqref{eq:exact-window-event}, without an added restriction on
$N_{m,R}$.
\end{corollary}
\begin{proof}
Here $q\le m+1=O(\log n)$. Substitution in
\eqref{eq:window-central-integral} gives its four polynomial margins
\eqref{eq:window-exponent-margins}, multiplied only by fixed powers
of $\log(2+n)$. Strict positivity absorbs those powers. For indicator
weights, invariance gives $\alpha_R=p_{n,R}$, not a product of the
individual probabilities. Divide the unconditional estimate by this
same probability. The strict inequalities with $\beta$ again absorb
the logarithmic factors. No induced mixing hypothesis is used.
\end{proof}

The admissible region is nonempty without a relation between the
complexity and return-tail constants. To exhibit one choice, enlarge
$C_{\mathrm w}$ to at least one and put
\begin{equation}\label{eq:nonempty-window-budget}
 \begin{split}
 d&=1/56,\qquad b=1/(56C_{\mathrm w}),\qquad
 a=\min\{b/2,c_{\mathrm t}b/(4a_{\mathrm t})\},\\
 \theta&=\min\{1/448,c_{\mathrm t}b/32\},\qquad
 \beta=\min\{1/448,c_{\mathrm t}b/16\},\qquad \kappa=0.
 \end{split}
\end{equation}
For every fixed $r$, these constants satisfy
$0<\theta<1/200$ and $0<\beta<\min_i\sigma_i$.
Indeed $C_{\mathrm w}b=d=1/56$ and
$a_{\mathrm t}a\le c_{\mathrm t}b/4$; hence the first three margins
are at least $11/448$, $12/448$, and $4/448$, respectively, and the
fourth is at least $5c_{\mathrm t}b/8$. This gives logarithmically
many observations with bounded individual variation and polynomially
rare exact events. The original $n^{1/200}$ band for the original
single-mark class remains Theorem C; \eqref{eq:window-exponent-margins}
states the cost of the enlarged class, rather than assigning it that
band without proof.

\subsection{Actual moments under the same event}
Set
\[
 \mathcal B=C(1+q\epsilon^{-1})e^{C_{\mathrm w}L},\qquad
 \mathcal E=\min\{1,\epsilon V+e^{a_{\mathrm t}m-c_{\mathrm t}L}\}.
\]
\begin{theorem}[Window-weighted actual moments]
\label{thm:window-weighted-moments}
With the same data, uniformly in $R,k$,
\begin{align}
 \left|\int W_R\circ(F_R^*)^k\,\frac{U_{n,R}}{\sqrt n}\dd\nu_R^*\right|
  &\le C\{n^{-1/6}+\mathcal B n^{-1/2}+\mathcal E^{1/2}\},
       \label{eq:window-first-moment}\\
 \left\|\int W_R\circ(F_R^*)^k\,
           \frac{U_{n,R}\otimes U_{n,R}}n\dd\nu_R^*
                           -\alpha_R D_R\right\|
  &\le C\{n^{-1/6}+\mathcal B n^{-1}+\mathcal E^{1/2}\}.
       \label{eq:window-second-moment}
\end{align}
For \eqref{eq:exact-window-event} with $p_{n,R}\ge p_0n^{-\beta}$,
the conditional mean tends to zero and the conditional covariance
tends to $D_R$ whenever
\begin{equation}\label{eq:window-moment-margins}
 \beta<\min\left\{\frac16,\frac12-C_{\mathrm w}b-d,
             \frac{d-\kappa}{2},
             \frac{c_{\mathrm t}b-a_{\mathrm t}a}{2}\right\}.
\end{equation}
The explicit choice \eqref{eq:nonempty-window-budget} satisfies these
inequalities as well.
\end{theorem}
\begin{proof}
Use $c_*^{-1}(W_{L,\epsilon})^0$ in
Theorem~\ref{thm:marked-quadratic-moments}. This gives the first two
terms of each bound with the approximant's exact mass.
The difference between the weights has modulus at most two, so
\eqref{eq:window-weight-L1} and invariance imply
\[
 \| (W_R-W_{L,\epsilon})\circ(F_R^*)^k\|_2
                               \le C\mathcal E^{1/2}.
\]
The actual second and fourth moments satisfy
$\|U_{n,R}/\sqrt n\|_2\le C$ and
$\|\,|U_{n,R}|^2/n\|_2\le C$ by Theorem D.
Cauchy--Schwarz therefore controls the two changes of weighted moment
by $C\mathcal E^{1/2}$. The change of mass is at most $C\mathcal E$
and $D_R$ is uniformly bounded. This proves
\eqref{eq:window-first-moment}--\eqref{eq:window-second-moment}.

For indicators, divide by the unchanged $p_{n,R}$. Under the
logarithmic choices, the four margins in
\eqref{eq:window-moment-margins} give convergence of the first moment
and of the second moment about the original center. Subtract the tensor
square of the conditional first moment to obtain the covariance.
For \eqref{eq:nonempty-window-budget}, these margins are at least
$1/6$, $13/28$, $1/112$, and $3c_{\mathrm t}b/8$, all strictly larger
than the displayed $\beta$.
\end{proof}

\subsection{The weighted raw inversion interface}
The new central estimate concerns the original window-weighted
measure
\[
 \mu^W_{n,R}=(J_{n,R})_*
                ((W_R\circ(F_R^*)^k)\nu_R^*).
\]
If the observation functions also belong to the bounded finite-record
subanalytic class of Theorem~\ref{thm:finite-count-raw-extraction},
that theorem supplies a finite-count extraction for this same measure.
This is an additional regularity requirement: arbitrary BV observations
are not thereby called subanalytic.

Here is the exact error decomposition which identifies the remaining
weighted estimates. Fix $0<\theta\le1/200$, set
$r_n=n^{-1/2+\theta}$, and choose a smooth cutoff $0\le\chi\le1$, equal
to one on the unit ball and zero outside the ball of radius two. Let
$\chi_{n,\theta}(\omega)=\chi(\omega/r_n)$ on a fixed neighborhood
of zero in $\T^3\times\R$, and let $K_{n,\theta}$ be its mixed
inverse Fourier kernel. Write $\mu^{W,L}=E^{W,L}+Q^{W,L}$ for the
finite extraction, $e^{W,L}$ for its extracted density, and
$T^{W,L}=\mu^W-\mu^{W,L}$. For a central point $y$ whose collision
count does not exceed $L$, the same algebra as in
\eqref{eq:count-localized-exact-inversion} gives
\begin{align}
 p^W(y)={}&(K_{n,\theta}*\mu^W)(y)
       +e^{W,L}(y)-(K_{n,\theta}*E^{W,L})(y)\notag\\
 &+\mathcal F^{-1}[(1-\chi_{n,\theta})\widehat Q^{W,L}](y)
       -(K_{n,\theta}*T^{W,L})(y).
       \label{eq:window-raw-exact-identity}
\end{align}
The inverse transform of the residual is legitimate by Theorem I.
The identity holds for the raw density representatives wherever those
are defined, and hence almost everywhere.

If $L=\lceil\lambda n\rceil$, $\lambda>c_*^{-1}+1$, every fixed
central window has the required count support for large $n$.
Since $|W_R|\le1$, the omitted high-count measure has total variation
at most one. The Poisson--Schwartz kernel estimate in
Section~\ref{sec:count-localized-inversion}, with $n^{1/200}$ replaced
by $n^\theta$, consequently gives
\[
 n^2\sup_{y\text{ in the central window}}
        |(K_{n,\theta}*T^{W,L})(y)|=O(n^{-P})
                    \quad\hbox{for every }P>0.
\]
Indeed its count labels are a distance proportional to $n$ away,
$r_n n=n^{1/2+\theta}$, and the normalized kernel prefactor is
$n^2r_n^4=n^{4\theta}$; any fixed polynomial loss is absorbed by a
sufficient Schwartz order.

Theorem~\ref{thm:window-central-integral} bounds the normalized error
in the first term of \eqref{eq:window-raw-exact-identity} by its
right-hand side, plus $C e^{-c n^{2\theta}}$ from the Gaussian tail.
The other terms require the actual weighted local-edge correction,
the complementary residual integral, and the growth of the finite
second-derivative sum. For an indicator window these requirements
must hold after division by the same $p_{n,R}$.
Thus the central and moment inputs now include a growing multiple-time
class. No weighted tail or exact physical-event replacement is inferred
from them; the latter still concerns a different event from
\eqref{eq:exact-window-event}.
